% CNNA architectural identity; conceptual, not an additional node theorem.
\section*{What CNNA means: ComplemeNt Net Architecture}
\addcontentsline{toc}{section}{What CNNA means: ComplemeNt Net Architecture}

The acronym \textbf{CNNA} expands to \textbf{ComplemeNt Net Architecture}.  The internal capital ``N'' is deliberate: the project is not merely about a growing tree or a collection of later algebras.  It studies one historically grown provenance \emph{net} through a recurring architecture of relative complements.

At a chosen growth stage, cut, scale, or projection, the same net is resolved schematically into
\[
  \text{realized or retained structure}
  \;\oplus\;
  \text{relative complement}
  \;\oplus\;
  \text{interface or handoff}.
\]
This display is architectural notation, not a claim that every later instance is a direct-sum decomposition.  The mathematical realization changes with the node: an unrealized continuation slot, a Schur interior, a cut-relative environment, or information discarded by a later observation channel can occupy the complementary role.

The architecture begins at the origin.  C001 fixes the empty baseline; C002 creates the first realized node; C003 and C018 distinguish the born prefix from admissible but not-yet-realized continuations; C004A selects the first open slot; and C013/C014 convert that slot into the first nontrivial weighted directed net.  Thus CNNA starts with growth itself, before Schur complements, quantum language, or operator algebras appear.

The same concept then persists through the derivation:
\[
\begin{aligned}
  &\text{open provenance slot}
  \longrightarrow \text{birth and directed relation},\\
  &\text{retained boundary / complementary interior}
  \longrightarrow \text{effective Schur/DtN response},\\
  &\text{immutable record / mutable live context}
  \longrightarrow \text{backreaction and memory},\\
  &\text{nested retained regions / nested complements}
  \longrightarrow \text{refinement and cocycle laws},\\
  &\text{provenance-complete state / reduced observation}
  \longrightarrow \text{later POVM, OQS, and AQFT specializations}.
\end{aligned}
\]

\begin{cnnaInfoBox}{Complement is a relative role, not an intrinsic node label}
CNNA does not assign every node one permanent opposite.  What counts as retained, complementary, or interface depends on the growth stage, cut, scale, and observation map.  A coordinate can be retained in one description and eliminated in another.  The name therefore denotes an architectural invariant of the construction, not an additional set-theoretic axiom.
\end{cnnaInfoBox}

\begin{cnnaInfoBox}{Formal status}
This section fixes terminology and the project-wide interpretation.  Its mathematical force remains exactly the force of the registered node contracts.  In particular, it does not promote the universal open-provenance hypothesis, the finite Legacy POVM, or later OQS/AQFT targets to proved results.
\end{cnnaInfoBox}
